\section{Semantic analysis}
\label{sec:semantic}

Semantic analysis checks the rules that the grammar cannot express. Some
checks happen when the program is compiled (static analysis) and others when
it runs (dynamic analysis).

\subsection{Static analysis}

Static analysis enforces semantic rules at compile time. Violating one of
these rules usually produces a compile error. Typical rules are the
visibility of variables, the checks on identifiers and other context
conditions, such as the number and types of the arguments of a call.

\paragraph{Validity and visibility rules.} When the compiler finds a variable,
as in \texttt{x = x + 1;}, it has to determine which declaration the name
refers to and whether that declaration is accessible at that point. In
\begin{lstlisting}
... {
  x = 5;
  x;
}
\end{lstlisting}
the question is whether a declaration of \texttt{x} is visible inside the
block. The lecture's harder case is a loop inside a loop with the same
variable:
\codelines{C}{c/nested_for.c}{6}{11}
In C and C++ each \texttt{for} opens a block, and the inner \texttt{int i}
declares a new variable that hides the outer one inside the inner loop. The
program prints $0, 1, \dots, 19$ ten times: 200 lines. The file
\texttt{nested\_for.c} checks the count. In Java the same code does not compile:
\texttt{variable i is already defined in method main(String[])}. Java forbids
a local variable from hiding another local variable of the same method. The
syntax is the same in both languages; the visibility rules are different.

\paragraph{Symbol tables chained by blocks.} The compiler keeps one table per
block and chains each table to the table of the enclosing block. In the
lecture's C example
\begin{lstlisting}[language=C]
int a; int b;
void p(); void q();
void p()
{ int a; int c; }
int c; int d;
void q()
{ int a; int d;
  void r()
  { int a; int c; }
}
\end{lstlisting}
the blocks are numbered in the order they open: block 1 is the global scope,
block 2 the body of \texttt{p}, block 3 the body of \texttt{q} and block 4 the
body of \texttt{r}. At the point inside block 4 the chain is the one of Figure~\ref{fig:chain}.
Block 2 is not in it: it was closed before block 3 started. A name is
searched from the innermost table outward, so \texttt{a} refers to the
declaration in block 4, \texttt{d} to the one in block 3 and \texttt{b} to
the global one.

\begin{figure}[h]
\centering
\begin{tikzpicture}[tab/.style={draw, rectangle split, rectangle split parts=#1, font=\small\ttfamily, minimum width=1.2cm}]
  \node[draw, font=\small] (b4) at (0,0) {4};
  \node[draw, font=\small] (b3) at (3,0) {3};
  \node[draw, font=\small] (b1) at (6,0) {1};
  \node[tab=2, below=0.4cm of b4] (t4) {a\nodepart{two}c};
  \node[tab=3, below=0.4cm of b3] (t3) {a\nodepart{two}d\nodepart{three}r};
  \node[tab=6, below=0.4cm of b1] (t1) {a\nodepart{two}b\nodepart{three}c\nodepart{four}d\nodepart{five}p\nodepart{six}q};
  \draw[flecha] (b4) -- (b3);
  \draw[flecha] (b3) -- (b1);
  \draw (b4) -- (t4);
  \draw (b3) -- (t3);
  \draw (b1) -- (t1);
\end{tikzpicture}
\caption{The chain of symbol tables inside block 4.}
\label{fig:chain}
\end{figure}

\subsection{Identification of identifiers}

\begin{definition}
The \textbf{identification of identifiers} links each applied occurrence of
an identifier (a use) to a defining occurrence (its declaration), or finds
that no such link exists or that more than one exists.
\end{definition}

Identification is based on six operations on the symbol table
(Table~\ref{tab:symtab-ops}).

\begin{table}[h]
\centering
\small
\begin{tabular}{l p{0.62\textwidth}}
  \toprule
  Operation & Meaning \\
  \midrule
  \texttt{create\_table()}      & creates an empty symbol table \\
  \texttt{enter\_block()}       & registers the opening of a new block \\
  \texttt{exit\_block()}        & resets the table to its state before the last \texttt{enter\_block} \\
  \texttt{enter\_id(id, ref)}   & enters \texttt{id} in the current block, with the reference \texttt{ref} to its declaration node \\
  \texttt{search\_id(id)}       & searches the defining occurrence of \texttt{id} in all the open blocks, innermost first; returns the reference or null \\
  \texttt{search\_block\_id(id)} & searches \texttt{id} only in the current block; returns the reference or null \\
  \bottomrule
\end{tabular}
\caption{Operations on the symbol table.}
\label{tab:symtab-ops}
\end{table}

In Rascal, with the table as a list of maps (the head is the current block):
\codelines{Rascal}{rascal/src/Scopes.rsc}{16}{27}

The lecture's language is
\codelines{Rascal}{rascal/src/Scopes.rsc}{6}{14}
and the word is \texttt{int a; \{ a = 4; int bar; \}}. Its parse tree is in
Figure~\ref{fig:scopes-tree}. In the attribute grammar of the lecture each
identifier leaf is a record $\langle \mathit{identifier}, \mathit{name}, \mathit{ref},
\mathit{succs}\rangle$ with the concrete name, a reference to its declaration and
the list of successors, and each interior node is $\langle \mathit{succs}, \mathit{rhs}\rangle$,
its children and the right side of the rule that produced it. The lexer
gives the identifier tokens their position: \texttt{<Id, bar, (l,c)>}.

\begin{figure}[h]
\centering
\begin{forest}
  for tree={font=\small\ttfamily, s sep=2mm, l sep=4mm},
  [Block
    [Decl [Type [int]] [id\textsubscript{a}] [;]]
    [Block
      [Stat
        [\{]
        [Block
          [Stat [id\textsubscript{a}] [=] [E [4]] [;]]
          [Block [Decl [Type [int]] [id\textsubscript{bar}] [;]] [Block [$\lambda$]]]
        ]
        [\}]
      ]
      [Block [$\lambda$]]
    ]
  ]
\end{forest}
\caption{Parse tree of \texttt{int a; \{ a = 4; int bar; \}}.}
\label{fig:scopes-tree}
\end{figure}

Each node has a procedure \texttt{process()} that does its part of the
identification, according to the rule of the node:
\begin{lstlisting}
process() {
  Decl ref;
  switch (rhs) {
    case "Type id;":
      ref := table.search_block_id(succs[1].id);
      if (ref != null) error();
      else table.enter_id(succs[1].id, this);
      return;
    case "id = E;":
      ref := table.search_id(succs[0].id);
      if (ref == null) error();
      else succs[0].ref := ref;
      return;
    case "{ Block }":
      table.enter_block();
      succs[1].process();
      table.exit_block();
      return;
  }
}
\end{lstlisting}
A declaration uses \texttt{search\_block\_id}: a name may be declared again in
an inner block (it hides the outer one), but not twice in the same block. A
use must use \texttt{search\_id}, because the declaration may be in any
enclosing block. On slide 37 the case \texttt{"id=E;"} uses
\texttt{search\_block\_id}; with that version the use of \texttt{a} inside the
block finds nothing, because \texttt{a} is declared in the outer block, and the
correct program is reported as wrong. Table~\ref{tab:process} follows the
corrected version on the lecture's word.

\begin{table}[h]
\centering
\small
\begin{tabular}{l l l}
  \toprule
  Node processed & Operation & Table after (current block first) \\
  \midrule
  start              & \texttt{create\_table()}                 & [ \{\} ] \\
  \texttt{int a;}    & \texttt{search\_block\_id(a)} = null, \texttt{enter\_id(a, 1)} & [ \{a$\mapsto$1\} ] \\
  \texttt{\{}        & \texttt{enter\_block()}                  & [ \{\}, \{a$\mapsto$1\} ] \\
  \texttt{a = 4;}    & \texttt{search\_id(a)} = 1, so \texttt{a.ref := 1} & [ \{\}, \{a$\mapsto$1\} ] \\
  \texttt{int bar;}  & \texttt{search\_block\_id(bar)} = null, \texttt{enter\_id(bar, 2)} & [ \{bar$\mapsto$2\}, \{a$\mapsto$1\} ] \\
  \texttt{\}}        & \texttt{exit\_block()}                   & [ \{a$\mapsto$1\} ] \\
  \bottomrule
\end{tabular}
\caption{Identification on \texttt{int a; \{ a = 4; int bar; \}}.}
\label{tab:process}
\end{table}

In Rascal, \texttt{process} walks the items of each block. The option
\texttt{slideVersion} switches the search of uses to \texttt{search\_block\_id},
and the test \texttt{slideError} shows the false error it produces:
\codelines{Rascal}{rascal/src/Scopes.rsc}{31}{67}

\paragraph{Decorated ASTs.} This paragraph goes beyond the slides. The result
of semantic analysis is the same tree with more information: the \texttt{ref}
of each use, the type of each expression, the kind of each operator (integer
or floating-point addition). Later phases read these annotations instead of
computing them again.

\subsection{Dynamic analysis}

Some rules can only be checked while the program runs: division by zero,
array indices out of bounds. Dynamic analysis is code that the compiler
generates, or that the programmer writes, to enforce them at run time. The
lecture names three forms.

\paragraph{Assertions.} A condition checked at a point of the program:
\codelines{C}{c/division.c}{4}{7}
If the condition is false the program stops with a message that says where.

\paragraph{Invariants.} A condition that must hold every time the program
passes through a point, for example at the start of each iteration of a
loop. For a loop that adds the first $n$ numbers, $s = 0 + 1 + \dots + (i - 1)$
holds at the start of every iteration.

\paragraph{Pre- and postconditions.} Conditions that must hold before and after
a piece of code. In the lecture's example
\[
\{x \le 10\}\ \ \texttt{while (x < 10) x = x + 1;}\ \ \{x = 10\}
\]
if $x \le 10$ before the loop, then $x = 10$ after it. The semantics block
returns to this notation (Hoare triples), where these conditions are proven
instead of checked.

\subsection{More static analyses}

This subsection goes beyond the slides.

\paragraph{Kinds of semantic errors.} Besides undeclared names and names
declared twice in a block, which the identification above detects, static
analysis reports arity mismatches (a call with a different number of
arguments than the declaration) and type mismatches (\texttt{3 + true}). The
lookup of a name can be written as two rules over the symbol table $\Gamma$:
\begin{mathpar}
\inferrule*[right=Found]{(x : \tau) \in \Gamma}{\Gamma \vdash x \Rightarrow \tau}
\and
\inferrule*[right=Undeclared]{x \notin \mathrm{dom}(\Gamma)}{\Gamma \vdash x \Rightarrow \textsc{undeclared}}
\end{mathpar}
When $\Gamma$ is a chain $\Gamma_{\mathit{inner}} :: \Gamma_{\mathit{outer}}$, the search
looks at $\Gamma_{\mathit{inner}}$ first, which gives shadowing.

\paragraph{Data-flow analysis.} Other analyses follow how values move through
the program. A definition (an assignment) \emph{reaches} a point if some path
goes from the definition to the point without reassigning the variable. In
\begin{lstlisting}
x = 1;
y = x;
x = 2;
print(x);
\end{lstlisting}
\texttt{x = 1} reaches \texttt{y = x} but not \texttt{print(x)}. If no definition
of a variable reaches a use, the use reads an uninitialized variable. A
variable is \emph{live} at a point if its current value may be read later
before being overwritten; a variable that is never live after an assignment
was assigned for nothing. Both analyses propagate facts along the
control-flow graph until nothing changes, like FIRST and FOLLOW.

\subsection{Exercises}

\begin{exercise}
What does the nested \texttt{for} of this section print in C? What happens in
Java? Change the Java version so that it prints the same as the C one.
\end{exercise}

\paragraph{Solution.} In C, the numbers 0 to 19, ten times. Java rejects the
program because the inner \texttt{i} would hide a local variable of the same
method. Renaming the inner variable to \texttt{j} and printing \texttt{j} gives
the same output.

\begin{exercise}
For the C program with blocks 1 to 4, give the chain of symbol tables at a
point inside block 2, and say which declaration each of \texttt{a}, \texttt{b},
\texttt{c} and \texttt{d} refers to there.
\end{exercise}

\paragraph{Solution.} The chain is block 2 $\to$ block 1. At that point block 1
only contains \texttt{a}, \texttt{b}, \texttt{p} and \texttt{q}, because
\texttt{c} and \texttt{d} are declared later in the text. So \texttt{a} and
\texttt{c} refer to block 2, \texttt{b} to block 1, and \texttt{d} is not
declared yet: using it there is an error.

\begin{exercise}
Run \texttt{process()} on \texttt{int a; \{ int a; a = 1; \} a = 2; b = a;}. Which
declaration does each use of \texttt{a} refer to, and which errors are
reported?
\end{exercise}

\paragraph{Solution.} The inner \texttt{int a;} is allowed: \texttt{search\_block\_id}
only looks at the inner block. The use \texttt{a = 1} finds declaration 2. After
\texttt{exit\_block()}, \texttt{a = 2} finds declaration 1. In \texttt{b = a},
\texttt{search\_id(b)} returns null and an error is reported; \texttt{a} refers to
declaration 1.

\begin{exercise}
In the program \texttt{x = 1; y = 2; if (y > 0) x = 3; print(x);}, which
definitions of \texttt{x} reach \texttt{print(x)}?
\end{exercise}

\paragraph{Solution.} Both. If the condition is false, \texttt{x = 1} reaches the
\texttt{print}; if it is true, \texttt{x = 3} does. A static analysis does not
know the value of \texttt{y} in general, so it keeps both.
